% !PRINTED 212
\weberpage{212}
%<<CONT
oder
\[
\begin{pmatrix} \alpha\,\varepsilon^r, & \beta \\ \gamma\,\varepsilon^r, & \delta \end{pmatrix} = \begin{pmatrix} \varepsilon^s\alpha, & \varepsilon^s\beta \\ \gamma, & \delta \end{pmatrix}.
\]

Wären $\beta$ und $\gamma = 0$, so wäre $\varphi$ selbst von der Form $\varepsilon x$,
und $\varepsilon$ eine $2m^{\text{te}}$ Einheitswurzel, und $D_m$ wäre also eine cyklische
Gruppe und keine Diëdergruppe. Ist aber $\beta$ oder $\gamma$ von Null
verschieden, so folgt $\varepsilon^s = \varepsilon^{-r}$ und $\alpha = \delta = 0$, also
\[
\varphi = \begin{pmatrix} 0, & \beta \\ \gamma, & 0 \end{pmatrix},
\]
was durch eine multiplicative Transformation auf die Form
$\begin{pmatrix} 0, & 1 \\ 1, & 0 \end{pmatrix}$ gebracht werden kann.

Unter den Diëdergruppen ist die Gruppe $D_2$ besonders
hervorzuheben, in der drei Systeme von je zwei zweizähligen
Polen vorhanden sind. Sie besteht aus den vier Substitutionen
\[
\begin{pmatrix} 1, & 0 \\ 0, & 1 \end{pmatrix},\quad
\begin{pmatrix} -1, & 0 \\ 0, & 1 \end{pmatrix},\quad
\begin{pmatrix} 0, & 1 \\ 1, & 0 \end{pmatrix},\quad
\begin{pmatrix} 0, & 1 \\ -1, & 0 \end{pmatrix}
\]
und wird die \gesperrt{Vierergruppe} genannt.

\begin{center}
\S.\,56. \\[6pt]
\gesperrt{Die Tetraëdergruppe.}
\end{center}

Bei der Tetraëdergruppe haben wir nach \S.\,52, III. ein
System von sechs conjugirten zweizähligen Polen. Wir können
also nach \S.\,53 annehmen, dass in der Gruppe die beiden Substitutionen

\[
\Theta(x) = -x, \quad \psi(x) = \frac{1}{x}
\]
vorkommen, und folglich auch
\[
\psi_1(x) = \psi\,\Theta(x) = -\frac{1}{x}.
\]

Die beiden Substitutionen $\psi$ und $\psi_1$ geben die Pole $\pm 1$
und $\pm i$, und diese müssen, da $\psi$ und $\psi_1$ vom $2^{\text{ten}}$ Grade sind
und in $G$ überhaupt nur zwei- oder dreizählige Pole vorkommen,
zu den zweizähligen gehören (\S.\,52, 3.). Die Substitution $\Theta(x)$
giebt die zweizähligen Pole $0, \infty$. Demnach lautet die Gleichung,
von der die zweizähligen Pole abhängen, $x(x^4-1) = 0$ und wir
erhalten die erste Grundform $6^{\text{ten}}$ Grades:
\weq{1}{f(x_1, x_2) = x_1 x_2 (x_1^4 - x_2^4).}
